\section{Part IV: The Inverse Problem - Data-Driven Approaches}
\label{sec:part4}
\lecture{01}{10}{2026}

% Mirrors slides.qmd "Part IV" outline. One subsection per outline item.

Parts~\partref{1}--\partref{3} assumed we know the governing equation
and want its solution. Part~\partref{4} inverts that assumption. What we
have instead are \textbf{snapshots} $\{u(t_k)\}$ --- from an experiment, from
DNS output, or from sensors. What we lack is the operator, the closure, and
sometimes the equation itself.

This is the regime of most real measurement data, and of every system too
complex to write down. The claim of this part is that the same compressed format
serves here too, and that it does so without giving up interpretability.

\subsection{Prediction of dynamics using time-series}
\label{subsec:prediction-time-series}
% Keywords: forecasting from snapshots, learning vs. prediction, snapshot
% matrices, curse of dimensionality in data-driven methods, why compressed
% formats help.

\paragraph{Two tasks, kept separate.}
\margintext{Conflating these two is the most common error in the data-driven
literature. A method can be excellent at one and useless at the other.}
\begin{itemize}
  \item \textbf{Prediction.} Given the past, forecast $u(t_{K+k})$. No
    interpretable model is required; only the forecast is judged.
  \item \textbf{Learning, or identification.} Recover the operator or the
    governing equations. Here interpretability is the entire point.
  \item Neural surrogates are the clean illustration: they predict well and
    explain nothing.
\end{itemize}

The classical toolkit spanning both tasks: POD and PCA, DMD and Koopman
operators, SINDy, operator learning (DeepONet, FNO), and data assimilation.

\paragraph{The same curse, again.}
\begin{itemize}
  \item Snapshot matrices are $N_x \times N_t$ with $N_x = 2^n$, so even taking
    the SVD is prohibitive --- the wall of
    \cref{subsec:numerical-methods-intro} in a new costume.
  \item The data brings its own problems on top: sparse and noisy measurements,
    short time series, non-stationarity, and no guarantee that a closed model
    exists at all.
  \item \textbf{The tensor-network promise:} compress \emph{both} space and time,
    perform the linear algebra inside the compressed format, and keep the modes
    interpretable.
\end{itemize}

\begin{remark}
  A useful practical bonus: a model obtained in MPS format plugs straight back
  into the Part~\partref{3} solver. Same data structure, no conversion
  layer, no lossy hand-off between the learned model and the simulator.
\end{remark}

\subsection{Dynamic Mode Decomposition and Space-Time Encoding}
\label{subsec:dmd-space-time}
% Keywords: DMD algorithm, Koopman operator viewpoint, snapshot matrices,
% space-time tensorization, QTT encoding of spatiotemporal data, connection
% to MPS-DMD.

\paragraph{DMD in one line.}
Assume the dynamics are approximately linear in the state,
\begin{equation}
  u_{k+1} \approx A\, u_k ,
  \label{eq:dmd-linear}
\end{equation}
find the best such $A$ from the data, and read off its eigenpairs. Each
eigenvalue gives a \textbf{growth rate and a frequency}; each eigenvector gives a
\textbf{spatial mode}. That is the whole appeal: the output is a set of
individually interpretable physical structures.

\paragraph{The standard recipe.}
\begin{itemize}
  \item Build snapshot matrices $X = [u_0 \cdots u_{T-1}]$ and
    $X' = [u_1 \cdots u_T]$.
  \item Take the SVD $X = U\Sigma V^{\dagger}$ and project onto the leading
    modes,
    \begin{equation}
      \tilde{A} = U^{\dagger} X' V \Sigma^{-1}.
      \label{eq:dmd-projected}
    \end{equation}
  \item Eigendecompose the small matrix $\tilde{A}$; the DMD modes are
    $\Phi = UW$.
\end{itemize}

\begin{remark}[Koopman view]
  Equation~\eqref{eq:dmd-projected} is a finite-rank approximation of the linear,
  infinite-dimensional Koopman operator acting on observables. That is why a
  linear method can describe nonlinear dynamics without contradiction.
\end{remark}

\paragraph{Space and time in a single MPS.}
\margintext{Our 2025 space--time paper. The conceptual move is small and the
consequence is large: once time is just more legs, there is no time loop to
parallelize or to accumulate error along.}
\begin{itemize}
  \item The idea~\cite{peddinti2025spacetime}: binarize \textbf{both} $x$ and
    $t$, into one MPS of $\log_2 N_x + \log_2 N_t$ tensors.
  \item Time stops being a loop and becomes simply more legs on the same chain.
  \item \textbf{Ordering matters here too.} Space-then-time works generally;
    time-then-space converged better for the nonlinear Schr\"odinger case,
    giving smoother vectorized solutions.
  \item The conceptual payoff: the \textbf{temporal} bond dimension becomes a
    direct measure of temporal complexity. Smooth dynamics means small $\chi$.
\end{itemize}

\paragraph{The all-at-once solve.}
\begin{itemize}
  \item Assemble the entire space--time linear system as a single MPO and solve
    it \textbf{once} with DMRG --- there is no time-stepping loop at all.
  \item Nonlinear terms are handled by Picard iteration, with linear
    convergence; Newton would give quadratic convergence at higher cost per
    iteration.
\end{itemize}

\noindent
Benchmarks on a $1024 \times 1024$ space--time grid:

\begin{center}
  \begin{tabular}{lll}
    \toprule
    \textsc{Equation} & $\chi$ & \textsc{Compression} \\
    \midrule
    Heat, wave                    & $2$--$22$ & $> 99\%$ \\
    Burgers, nonlinear diffusion  & $10$      & $> 99\%$, residual $\sim\!10^{-3}$ \\
    \bottomrule
  \end{tabular}
\end{center}

\paragraph{MPS-DMD.}
Standard DMD needs an SVD of an $N_x \times N_t$ matrix. That is the bottleneck
--- and, once the data is already a space--time MPS, it is unnecessary.

\begin{itemize}
  \item \textbf{MPS-DMD}~\cite{peddinti2025spacetime}: run the entire DMD
    pipeline \emph{inside} the space--time MPS, never forming the snapshot
    matrix.
  \item The key observation is that the MPS already contains the SVD. It only
    has to be gauged into the right form --- see \cref{subsec:mps}.
  \item Same modes, same eigenvalues, same predictions, at exponentially lower
    cost.
\end{itemize}

\noindent
The algorithm is three steps:

\begin{enumerate}
  \item \textbf{Gauge} into mixed canonical form at the space/time bond. The
    Schmidt decomposition there \emph{is} $U\Sigma V^{\dagger}$, so the SVD comes
    for free.
  \item \textbf{Project} the evolution operator onto the POD basis,
    $\tilde{A} = U^{\dagger} A U$: a small $\chi \times \chi$ matrix.
  \item \textbf{Eigendecompose} $\tilde{A} = W\Lambda W^{-1}$, take modes
    $\Phi = UW$, and predict by local contractions alone,
    \begin{equation}
      u_{T+k} = A^k u_T = U W \Lambda^k W^{-1} U^{\dagger} u_T .
      \label{eq:mps-dmd-predict}
    \end{equation}
\end{enumerate}

For externally supplied data --- where no space--time solve produced the MPS ---
stack the snapshot MPSs sequentially, truncating after each. The cost is linear
in the number of timesteps.

\paragraph{Results.}
\margintext{The headline is the scaling, not the compression ratio: runtime
logarithmic in \emph{both} $N_x$ and $N_t$, where prior tensor DMD variants are
linear in $N_t$.}
\begin{center}
  \begin{tabular}{@{}lp{0.62\textwidth}@{}}
    \toprule
    \textsc{Case} & \textsc{Setup and outcome} \\
    \midrule
    Burgers
      & $\chi = 10$ on $1024^2$; accurate $1000$ steps ahead, at $99\%$
        compression \\
    2D cylinder wake
      & $256\times128\times1024$ as a 25-bit MPS at $\chi = 200$; $97.42\%$
        compression, with 100 modes capturing the vortex street \\
    \bottomrule
  \end{tabular}
\end{center}

\subsection{Recovering Physics from Data}
\label{subsec:recovering-physics}
% Keywords: learning vs. prediction tasks, system identification, data
% assimilation, sparse/noisy measurements, governing-equation discovery.

DMD recovers an operator. The more ambitious task is to recover the
\emph{equation} --- and the difference matters, because an operator is tied to
the regime it was fitted in whereas an equation is not.

\begin{itemize}
  \item \textbf{System identification.} Fit a parametrized operator or a
    dictionary of candidate terms, and select among them. SINDy-style sparse
    regression is the canonical example; TT-based variants keep the dictionary
    manageable at high resolution.
  \item \textbf{Data assimilation.} Combine a known but imperfect model with
    sparse observations. Here the tensor network carries the state, and the
    model is the Part~\partref{3} solver.
  \item \textbf{Sparse and noisy measurements} are the normal case, not the
    exception. Recall from \cref{subsec:finding-qtt} that cross interpolation
    will fit noise faithfully if you let it.
  \item \textbf{Governing-equation discovery} is the hardest of these, and the
    one where interpretability is not a bonus but the deliverable.
\end{itemize}

\begin{remark}
  The truncation that makes these methods affordable is also, viewed
  differently, a regularizer: discarding small singular values discards
  low-amplitude structure, some of which is noise. Whether it denoises or
  faithfully reproduces the noise depends on where the noise sits in the
  spectrum, and that is not something you get to assume.
\end{remark}

\subsection{Tensor Networks for Learning Dynamical Laws}
\label{subsec:tn-learning-dynamics}
% Keywords: TN-based surrogate models, operator learning in TT format,
% training/optimization strategies, generalization and extrapolation in
% time, comparison to neural-network surrogates.

\begin{itemize}
  \item \textbf{TT surrogates and operator learning.} Represent the solution
    operator itself as an MPO and fit it to data. The training problem is again
    a sweeping optimization, so the machinery of
    \cref{subsec:linear-solvers} carries over directly.
  \item \textbf{Tensor-network closures.} Learn only the unresolved part --- the
    closure term --- and leave the resolved dynamics to the solver. This is the
    hybrid route, and the one where the physics prior does the most work.
  \item \textbf{Training strategies.} Sweeping optimization with adaptive ranks,
    warm starts across parameter values, and rank as the primary regularization
    knob.
  \item \textbf{Generalization and extrapolation.} Extrapolating in time is
    where surrogate models usually fail. A linear-in-observables model of the
    DMD kind extrapolates in a way you can reason about; a neural surrogate
    typically does not.
\end{itemize}

\begin{remark}[The comparison that matters]
  Against neural-network surrogates the tensor-network case does not rest on raw
  accuracy. It rests on the fact that modes, frequencies and growth rates come
  out of the method directly, rather than from a post-hoc probe --- and on the
  fact that rank gives a single, monitorable complexity knob.
\end{remark}

\subsection{Outlook}
\label{subsec:part4-outlook}
% Keywords: open challenges in data-driven TN methods, hybrid
% physics+data approaches, scalability, uncertainty quantification.

\paragraph{Extending the method.}
\margintext{The $(1+1)$D restriction is the most immediate limitation of the
space--time work, and lifting it brings back the index-ordering questions of
\cref{subsec:qtt-functions-operators}.}
\begin{itemize}
  \item \textbf{Beyond $(1+1)$D.} The space--time analysis so far is
    one-dimensional in space. Two and three dimensions plus time --- and the
    index ordering that goes with them --- is open.
  \item \textbf{Nonlinear and Koopman extensions.} Extended DMD with TT
    dictionaries, and kernel DMD in compressed form.
  \item \textbf{Noise and sparsity.} How TT truncation interacts with
    measurement noise is not settled; it can denoise, or it can fit the noise.
  \item \textbf{Online and streaming.} Update the MPS as data arrives, without
    recompressing from scratch.
\end{itemize}

\paragraph{Practice and positioning.}
\begin{itemize}
  \item \textbf{Hybrid physics and data.} Use the Part~\partref{3} solver as
    a prior and learn only the closure or the residual in TT form.
  \item \textbf{Uncertainty quantification.} Error bars on tensor-network
    predictions are essentially absent today. This is a serious gap for any
    engineering use.
  \item \textbf{Interpretability is the selling point.}
  \item \textbf{Software and reproducibility.} Shared TT/MPS tooling, standard
    benchmarks, and open datasets.
\end{itemize}
